\section{Overview}
\label{m:sec:overview}

This chapter sets out the mathematical framework used throughout the thesis: Markovian
population models, discrete or continuous in time, in which counts may be small enough
that extinction, establishment, and absorption are random events rather than the
deterministic outcomes of a rate equation. When the mean trajectory answers the wrong
question, one needs a distribution, and the apparatus below is chosen for that purpose.

The chapter is organised as objects first, methods second.

\Cref{m:sec:markov} surveys the process classes that appear later. In discrete time the
main examples are simple random walks and Galton--Watson branching (the latter only
briefly here). In continuous time the building blocks are exponential clocks and the
chains assembled from them: Poisson processes, birth--death processes, and
birth--death--catastrophe processes. Time-inhomogeneous models and systems that couple
ordinary differential equations to continuous-time Markov chains complete the inventory.
Notation for generators, competing clocks, and probability generating functions is fixed
along the way.

\Cref{m:sec:methods} collects the analytic methods applied to those processes.
Discrete-time devices---first-step analysis and functional iteration of generating
functions---are treated briefly; a fuller development of the quasi-stationary amplitude
$A(p)$ and of Koenigs-function methods is deferred to \ChA. Continuous-time methods are
developed with linear birth--death processes as the running example: backward equations,
means and higher moments, hitting and extinction probabilities, and conditional means,
including the continuous-time constant $\Ac(p)$ and a light definition of the discrete
constant $A(p)$. The method of characteristics is presented with one worked example;
extended absorption models are recorded in the appendix.

Two modelling threads run through later chapters and are introduced only as far as the
toolkit requires. The first is \emph{conditioning on survival}: a subcritical branching
process dies out almost surely, yet conditioned on non-extinction its population settles
to a quasi-stationary law whose mean is $1/A(p)$. The second is the \emph{near-critical
regime}, in which lifetimes become long and numerical schemes that iterate to stationarity
become expensive. Deep analysis of $A(p)$ itself---product and series representations,
bounds, closed-form search, and the Koenigs obstruction---belongs to \ChA, not to the
present chapter.

Open questions that sit with this toolkit rather than with the arithmetic of $A(p)$ include
the existence of a quasi-stationary distribution when one conditions on neither extinction
nor rupture, and the tightness of Jensen-type bounds for discrete catastrophe survival.
Those are left for later modelling chapters.
